\chapter{Automata: Two-Loop, Second-Order Systems (2-OS)}\label{lect3}
\localtoc

In this chapter, we will close a second feedback loop in digital systems. This loop will increase the autonomy of the system. If in 0-OS, because we had no loop, the output of the combinational systems is a combination that derives strictly from the current value applied to the input, then in 1-OS a partial autonomy was obtained: the autonomy of the state that could be maintained outside the range of during which the signal that determined it acted. The second loop, which defines 2-OS, will allow the autonomy of the behavior of the system in which it closes. That is, we will obtain the possibility of an evolution of the output even in the absence of a variation of the input signal.

The main targets in this chapter are the main components of a small and simple processing unit: {\bf {\em Control Automaton}}, {\bf {\em Program Counter (PC)}} and  {\bf {\em Registers with Arithmetic \& Logic Unit (RALU)}}.

%%%%%%%%%%%%%%%%%%%%%%%%%%%%%%%%%
\section{Historical Development}
%%%%%%%%%%%%%%%%%%%%%%%%%%%%%%%%%
Automata have a long history, evolving from ancient myths to modern computational models. Ancient myths featured self-operating machines, such as \textit{Talos}, a bronze automaton from Greek mythology \cite{bedini1964early}. The engineer Hero of Alexandria (10–70 AD) described pneumatic and hydraulic devices capable of motion \cite{hero1899pneumatics}. In China, Su Song (1020–1101 AD) created an elaborate water-powered astronomical clock with mechanical figurines \cite{needham1986science}.

During the Renaissance, Leonardo da Vinci (1452–1519) designed mechanical knights \cite{rosheim2006leonardo}. The 17th century saw Pascal's mechanical calculator \cite{pascal1665arithmetical}. Islamic engineers, such as Al-Jazari (1136–1206 AD), constructed sophisticated water clocks with programmable features \cite{aljazari1974ingenious}.

The Jacquard loom (1804) pioneered programmability \cite{essinger2004jacquard}. Charles Babbage (1791–1871) and Ada Lovelace envisioned the Analytical Engine \cite{menabrea1842sketch}. In the late 19th century, Hermann Hollerith developed the punched card system for data processing \cite{hollerith1890tabulating}.

Turing's work (1936) established the Turing Machine \cite{turing1936computable}, while von Neumann explored self-replicating automata \cite{von1951general}. Chomsky classified formal languages \cite{chomsky1956three}. In the mid-20th century, Shannon's work on Boolean circuits laid the foundation for digital logic \cite{shannon48}.

John Conway’s \textit{Game of Life} (1970) popularized cellular automata \cite{gardner1970mathematical}. Today, automata influence AI, linguistics, robotics, and formal verification \cite{hopcroft2006automata}. Quantum computing and bio-inspired automata are emerging fields extending automata theory’s applications \cite{nielsen2010quantum}.

%%%%%%%%%%%%%%%%%%%%%%
\section{Definitions}
%%%%%%%%%%%%%%%%%%%%%%

The typical circuit in 2-OS is the automaton. We will highlight two categories of automata: small \& complex finite automata and large \& simple functional automata.

%*******************************
\subsection{Generic Definition}
%*******************************
The following definition applies to all kinds of automata.

\begin{definition}
An automaton, A, is defined by the following 5-tuple:
$$A=(X, Y, Q, f, g)$$
where:
\begin{description}
\item[X]: the finite set of input variables \item[Y]: the finite
set of output variables \item[Q]: the set of state variables
\item[f]: the state transition function, described by $f:X \times
Q \rightarrow Q$ \item[g]: the output transition function,  with
one of the  following  definitions:
\begin{itemize}
    \item $g:X \times Q \rightarrow Y$ for the Mealy type automaton
    \item $g:Q \rightarrow Y$ for the Moore type automaton
    \item $g(q)=q$ for $Y \equiv Q$, where $q \in Q$ for the half-automaton, symbolized with $A_{1/2}$.
\end{itemize}
\end{description}
At each clock cycle, the state of the automaton  switches, and the
output takes the value according to the new state (and the current
input, in Mealy's approach).

$\diamond$
\end{definition}

A strict initial automaton is defined by:
$$A=(X, Y, Q, f, g; q_{0})$$
and has a special input, called {\bf reset}, used to set the
automaton to the initial state $q_{0}$. If the automaton is
initial, the input reset switches the automaton in one,
specially selected, initial state.

%*******************************
\subsection{Size vs. complexity}
%*******************************

The huge {\em size} of the actual circuits implemented on a single
chip imposes a more precise distinction between {\em simple}
circuits and {\em complex} circuits. When we can be integrated on a
single chip more than a Google ($10^{10}$) of components, 
the size of the circuits we design 
becomes less important than their complexity. We will
make a clear distinction between {\em size} and {\em
complexity}. We usually see imprecisely: ``the complexity of a computation is
given by the size of memory and by the CPU time". But, if we have
to design a circuit of 100 million transistors, it is very important
to distinguish between a circuit having a uniform structure and a randomly structured one. 
In the first case, the circuit can be
easily specified, easily described in an HDL, easily tested, and so on.
Otherwise, if the structure is completely random, without any
repetitive substructure inside, it can be described using only a
description having a similar dimension to the circuit size. When
the circuit is small, it is not a problem, but for billions of
components, the problem has no solution. Therefore, if the circuit
is very big, it is not enough to deal only with its size. The most
important becomes also the degree of uniformity of the circuit.
This degree of uniformity, the degree of order inside the circuit,
can be specified by its {\bf complexity}.

As a consequence, we distinguish more carefully the concept of
{\em size} and the concept of {\em complexity}. 

\begin{definition}
The {\bf size} of a digital circuit, $S_{digital \; circuit}$, is
given by the dimension of the physical resources used to implement

it.
$\diamond$
\end{definition}

In order to provide a numerical expression for size, we need a more
detailed definition that takes into account technological
aspects. In the '40s we counted electronic bulbs, in the '50s we
counted transistors, in the '60s we counted SSI\footnote{{\em
Small Size Integrated} circuits} and MSI\footnote{{\em Medium Size
Integrated} circuits} packages. In the '70s we started to use two
measures: sometimes the number of transistors (or the number of
2-input gates on the silicon die) and/or the silicon die
area. Thus, we adopt two numerical measures for size.

\begin{definition}
The {\bf gate size} of a digital circuit, $GS_{digital \;
circuit}$, is given by the total number of CMOS pairs of
transistors used for building the gates  used to implement it\footnote{Sometimes gate size
is expressed in the total number of 2-input gates necessary to
implement the circuit. We prefer to count CMOS pairs of
transistors (almost identical with the number of inputs) instead
of equivalent 2-input gates because it is the simplest.}.

$\diamond$
\end{definition}

This definition of gate size provides a partially accurate estimate of
the silicon area required to implement the circuit. The effects of
layout, fan-out, and speed are not captured by this definition.

\begin{definition}
The {\bf area size} of a digital circuit, $AS_{digital \;
circuit}$, is given by the dimension of the area of silicon used
to implement it.

$\diamond$
\end{definition}

The area size is useful to compute the price of the implementation
because when a circuit is produced, we pay for the number of
wafers. If the circuit has a big area, the number of the circuits
per wafer is small, and the yield is low\footnote{The same number
of errors make useless  a bigger area of the wafer containing
large circuits.}.

\begin{definition}\label{circCompl}
The {\bf algorithmic complexity} or simply the {\bf complexity} of a digital circuit, $C_{digital \; circuit}$, has the magnitude
order given by the minimal number of symbols needed to express its
definition.

$\diamond$
\end{definition}

Definition \ref{circCompl} is inspired by \cite{Chaitin77}'s definition for the
algorithmic complexity of a string of symbols. The
{\em algorithmic complexity} of a string is related to the
dimension of the smallest program that generates it.  Our $C_{digital
\; circuit}$ can be associated with the shortest unambiguous circuit description
in a certain HDL (in most cases, it is a behavioral description).

\begin{definition}
A {\bf simple circuit} is a circuit having a complexity
much smaller than   its size:
$$C_{simple \; circuit} << S_{simple \; circuit}.$$
Usually, the complexity of a simple circuit is constant:
$C_{simple \; circuit} \in O(1)$.

$\diamond$
\end{definition}

\begin{definition}
A {\bf complex circuit} is a circuit having the complexity  in
the same order of magnitude with its size:
$$C_{complex \; circuit} \sim S_{complex \; circuit}$$

$\diamond$
\end{definition}

\begin{exemplu}
The following Verilog program describes a {\em complex circuit},
because the size of its definition (the program) is
$$S_{def. \; of \; random\_circ} = k_{1} + k_{2} \times
S_{random\_circ} \in O(S_{random\_circ}).$$
{\em 
\includeverilogcode[caption={Example of a complex circuit.  File: A10\_randomCircuit.sv  (SystemVerilog)}, 
                    label={lst:A10_randomCircuit.sv}, 
                    firstline=1, 
                    firstnumber=1, 
                    lastline=30]
                    {\verilogSourceDir/A10_randomCircuit.sv}
}

$\diamond$
\end{exemplu}
Listing \ref{lst:A10_randomCircuit.sv} shows \texttt{A10\_randomCir} that defines a combinational circuit with five input signals: \( a, b, c, d, e \), and two output signals: \( f \) and \( g \). Additionally, it includes two internal logic signals, \( w1 \) and \( w2 \), which serve as intermediate nodes in the circuit.

The circuit begins by computing \( w1 \) using an AND gate, which takes \( a \) and \( b \) as inputs. Similarly, another AND gate computes \( w2 \) using inputs \( c \) and \( d \). 
The output \( f \) is determined by an OR gate, which takes \( w1 \) and \( c \) as inputs. Expanding the expression, we get:

\[
f = w1 + c = (a \cdot b) + c
\]

The second output, \( g \), is computed using another OR gate, which takes \( e \) and \( w2 \) as inputs. This results in the equation:

\[
g = e + w2 = e + (c \cdot d)
\]

This module implements a combinational logic circuit using basic AND and OR operations to determine the outputs \( f \) and \( g \) based on the provided inputs.


\begin{exemplu}
The following Verilog program describes a {\em simple circuit},
because the program that defines the circuit completely is the same
for any value of {\em n}.
{\em 
\includeverilogcode[caption={Example of a simple circuit.  File: A10\_orPrefixes.sv  (SystemVerilog)}, 
                    label={lst:A10_orPrefixes.sv}, 
                    firstline=1, 
                    firstnumber=1, 
                    lastline=30]
                    {\verilogSourceDir/A10_orPrefixes.sv}
}
The {\em prefixes of the OR} circuit consists of {\em n} $OR_{2}$ gates
connected in a very regular form. The definition is the same for
any value of {\em n}.\footnote{To be rigorous, when the
dimension of the input is specified, the
parameter $n$ is expressed using a string o symbols in $O(log \;
n)$. Typically, this aspect can be ignored.}

$\diamond$
\end{exemplu}
Listing \ref{lst:A10_orPrefixes.sv} shows \texttt{orPrefixes} that implements a parallel prefix computation using bitwise OR operations. The module is parameterized by \( n \), which defaults to 256, meaning it can operate on bit vectors of varying lengths.

The module has an output signal \texttt{out} and an input signal \texttt{in}, both of which are \( n \)-bit wide logic vectors. The purpose of this module is to compute the prefix OR of the input vector and store the results in the output vector.

The process is handled within an \texttt{always\_comb} block, ensuring that the computation occurs continuously as a purely combinational operation. The first element of the output, \texttt{out[0]}, is assigned directly from the first element of the input, \texttt{in[0]}, as shown in the equation:

\[
\text{out}[0] = \text{in}[0]
\]

A \texttt{for} loop iterates from index \( k = 1 \) to \( k = n-1 \), computing each subsequent element of the output vector as the bitwise OR of the corresponding input bit and the previously computed output bit. This results in the recurrence relation:

\[
\text{out}[k] = \text{in}[k] \ | \ \text{out}[k-1]
\]

Expanding this relation, each element in \texttt{out} accumulates the OR of all previous bits in \texttt{in}, forming a prefix OR operation:

\[
\text{out}[1] = \text{in}[1] \ | \ \text{in}[0]
\]

\[
\text{out}[2] = \text{in}[2] \ | \ \text{in}[1] \ | \ \text{in}[0]
\]

\[
\text{out}[3] = \text{in}[3] \ | \ \text{in}[2] \ | \ \text{in}[1] \ | \ \text{in}[0]
\]

This pattern continues until \texttt{out[n-1]}, which contains the OR of all bits in \texttt{in} from index \( 0 \) to \( n-1 \).

The module effectively computes a cumulative OR function over the input vector, making it useful in applications such as prefix sum computations, data compression, or bitmask operations where a running OR of previous elements is required. Since it uses a simple loop in an \texttt{always\_comb} block, it can be synthesized efficiently in hardware without requiring clocked registers.


\begin{verisum}{\em :

\hspace{-0.8cm} \rule[5mm]{16cm}{0.25mm} \vspace{-0.7cm}

\begin{description}
    \item[and, or, xor, not]: are reserved names for predefined circuits that can be instantiated under specific names
    \item[parameter]: used to define a local parameter
    \item[integer] k: defined a positive integer
    \item[for]: defined the well-known loop running under the index {\tt k}
\end{description}
\rule[5mm]{16cm}{0.25mm} }
\end{verisum}

Composing circuits generates not only the largest structures but also the
{\em deepest} ones. The depth of the circuit is related to the
associated propagation time.

\begin{definition}
The {\em depth} of a combinational circuit is equal to the total
number of serially connected {\em constant} input gates (usually 2-input gates) on the longest path from inputs
to the outputs of the circuit.

$\diamond$
\end{definition}



At the current technological level, the size has become less
important than the complexity because we can {\em produce}
circuits having an increasing number of components, but we can
{\em describe} only circuits having the range of complexity
limited by our mental capacity to deal efficiently with complex
representations. The first step to have a circuit is to express
what it must do in a behavioral description written in a certain
HDL. If this "definition" is too large, having the magnitude order
of a huge multi-billion-transistor circuit, we don't have the
possibility to write a testable program expressing our desire.

In the domain of circuit design, we passed long ago beyond the
stage of {\em minimizing} the number of gates in a few gates
circuit. Now, the most important thing in the
multi-billion-transistor circuit era, is the {\em ability to
describe} a circuit simple enough (because we can't
write huge programs), and sized large enough (because we can produce more circuits on
the same area). We must take into consideration that Moore's Law applies to size not
to complexity.

%**********************
\subsection{Taxonomy}
%**********************
Starting from the generic definition, Figure \ref{fin_aut} represents the generic structures of the automata used in digital design.

The structures found in current practice are:
\begin{description}
    \item[half automaton]: is an automaton whose output is identical to the internal state (the combinational circuit that transforms the state into the output is missing); the latency of the Y output compared to the X input is one clock cycle

\placedrawing[H]{figures/04_01_automata_types.lp}{{\bf Automata types.} {\small {\bf
a.} The structure of the half-automaton ($A_{1/2}$), the no-output function
automaton: the state is generated by the previous state and the
previous input. {\bf b.} The logic symbol of half-automaton. {\bf
c.} Immediate Mealy automaton: the output is generated by the
current state and the current input. {\bf d.} Immediate Moore
automaton: the output is generated by the current state. {\bf e.}
Delayed Mealy automaton: the output is generated by the previous
state and the previous input. {\bf f.} Delayed Moore automaton:
the output is generated by the previous state.}}{fin_aut}
    
    
    \item[Mealy immediate automaton]: is an automaton whose output is calculated by a combinational circuit depending on state and input; the latency of the Y output compared to the X input is zero
    \item[Moore immediate automaton]: is an automaton whose output is calculated by a combinational circuit depending only on the state; the latency of the Y output compared to the X input is one clock cycle
    \item[Mealy delayed automaton]: is a Mealy immediate automaton with the output delayed through a pipeline register; the latency of the Y output compared to the X input is one clock cycle
    \item[Moore delayed automaton]: is a Moore immediate automaton with the output delayed through a pipeline register; the latency of the Y output compared to the X input is two clock cycles.
\end{description}
Depending on the way the machine is integrated into the system we are designing, we will choose the most suitable version. Latency can sometimes be advantageous.




We will make another important distinction: that between complex automata and simple ones.

\begin{definition}
A complex automaton has a description proportional to the number of states it has.

$\diamond$
\end{definition}

\begin{definition}
A simple automaton has a description of constant size independent of the number of its states.

$\diamond$
\end{definition}

%%%%%%%%%%%%%%%%%%%%%%%%%%%%%%%%%%%%
\section{Finite (complex) Automata}
%%%%%%%%%%%%%%%%%%%%%%%%%%%%%%%%%%%%
The term "finite automata" does not imply a distinction from a possible category of "infinite automata." Rather, the adjective "finite" refers specifically to the number of states in the automaton, which remains bounded regardless of the length of the input it processes. A finite automaton is capable of handling an arbitrarily long, or even conceptually infinite, sequence of symbols presented at its input. Thus, the finiteness of the automaton pertains to the structure of its state space, not to any limitation on the length of the input string it can process. In this sense, the automaton's finiteness is contrasted with the metaphorical infinity of the input rather than an actual infinity of the automaton itself.

\begin{definition}
A finite automaton is characterized by $|Q| \in O(1)$, i.e., the size of the set Q is constant and independent of the length of the string of symbols it receives.

$\diamond$
\end{definition}

We will illustrate this concept using two representative cases: a recognition automaton and a control automaton.

%********************************
\subsection{Recognizing automata}
%*********************************
Theoretical linguist Noam Chomsky introduced the concept of generative grammar \cite{Chomsky56, Chomsky59, Chomsky63}, a concept used in modern computer science to develop compiler theory, a theory used to translate programs written in a high-level language into the machine language used at the lowest level in a computing machine \cite{Hopcroft2001}. The first step in translation is to recognize the sequence to be translated as belonging to a particular language. The simplest possible example concerns the languages defined in the set $\{0,1\}^*$, that of all strings formed with elements from the set $\{0,1\}$. We will consider, in the following example, a language from the category of regular languages, in the sense of the classification made by N. Chomsky. Regular languages, or type 3, are the simplest languages in the sense that they are generated with the most restrictions: the string is formed by increasing each step independently of the context at one end only. For these languages, a finite automaton (a two-loop digital system) can be used to recognize or generate.

\begin{exemplu}\label{regLang}
The binary inputs to the language $L = \{1^{n}0^{m} | n,m \geq 1\} \in \{0,1\}^*$,
are recognized by a finite half-automaton.
Let's define and design it. The transition diagram defining the
behavior of the half-automaton is presented in Figure \ref{rega},
where the state set $Q$ contains:

\placedrawing[H]{figures/04_02_transition_diagram.lp}{{\bf Transition diagram.} {\small
The transition diagram for the half-automaton, which recognizes
strings of form $1^{n}0^{m}$, for $n \geq 1$ and $m \geq 1$. Each
circle represents a state, and each (marked) arrow represents a
(conditioned) transition.}}{rega}

\begin{itemize}
    \item $q_{0}$ - is the {\em initial} state in which:
        \begin{itemize}
            \item if 1 is received, the output is {\tt rec1} and the automaton switches in the state $q_{1}$
            \item if 0 is received the automaton generate on output {\tt fail} and switches in $q_{3}$
        \end{itemize}
    \item $q_{1}$ - in this state at least one 1 was received and
        \begin{itemize}
            \item if 1 is received, the output is {\tt rec1} and the state remains the same
            \item if 0 is received, the output is {\tt rec0} the the state becomes $q_{2}$
        \end{itemize}
    \item $q_{2}$ - this state acknowledges a well formed string:
        \begin{itemize}
            \item if 1 is received, the output is {\tt fail} and the state switches in $q_{3}$
            \item if 0 is received, the output is {\tt rec0} and the state remains the same
        \end{itemize}
    \item $q_{3}$ - is the  error state because  an incorrect string was; the state remains unconditionally the same until a new {\tt reset} is aplied.
    received.
\end{itemize}



The first step in implementing the structure of the just-defined
automaton is to {\bf assign binary codes} to each state.
In this stage, we have absolute freedom. Any assignment can be
used. The only difference will be in the resulting structure but
not in the resulting behavior.

\begin{table}[H]
  \begin{center}
    \caption{Transition state table.}
    \label{table1}
    \begin{tabular}{|c|c|c||c|c|c|c|}
      \hline
     $Q_1$ & $Q_0$& $X$& $Q^+_1$& $Q^+_0$ & $Y_1$ & $Y_0$\\
     \hline \hline
   0&       0&     0&        0&       0 & 1  & 1 \\ \hline
   0&       0&     1&        0&       0 & 1  & 1 \\ \hline
   0&       1&     0&        0&       1 & 0  & 1 \\ \hline
   0&       1&     1&        0&       0 & 1  & 1 \\ \hline
   1&       0&     0&        0&       0 & 1  & 1 \\ \hline
   1&       0&     1&        1&       1 & 1  & 0 \\ \hline
   1&       1&     0&        0&       1 & 0  & 1 \\ \hline
   1&       1&     1&        1&       1 & 1  & 0 \\ \hline
    \end{tabular}
  \end{center}
\end{table}

\placedrawing[h]{figures/04_03_4state_finite_automaton.lp}{{\bf A 4-state finite
automaton.} {\small The structure of the finite
automaton used to recognize binary string belonging to the
$1^{n}0^{m}$ set of strings.}}{aut1}

Let be the codes, symbolized by $Q_1$ and $Q_2$, assigned in square brackets
in Figure \ref{rega}. For the outputs, the symbols $Y_1$ and $Y_2$ are used with the following meanings:
\begin{verbatim}
        rec0: 01
        rec1: 10
        fail: 11
\end{verbatim}
The input signal is symbolized by $X$.

The results of the transition table, for the functions $f$ and $g$, are presented in
Table \ref{table1}. The resulting transition functions of the resulting immediate Mealy automaton are:
$$Q_{1}^{+} = Q_{1} \cdot X = ((Q_{1} \cdot X)')'$$
$$Q_{0}^{+} = Q_{1} \cdot X + Q_{0} \cdot X' = ((Q_{1} \cdot X)' \cdot (Q_{0} \cdot X'))'$$
$$Y_1 = (Q_0 \cdot X)'$$
$$Y_0 = Q_1 \cdot X' = (Q_1' + X)'$$

The circuit is represented in Figure \ref{aut1} in a version using inverted gated
only. The 2-bit state register is designed with 2-D flip-flops. The
{\tt reset}  input is applied on the {\em set} input of D-FF1 and
on the {\em reset} input of D-FF0.

The recognition process begins with the application of the {\tt reset} signal through which the automaton goes to the initial state, $q_0$, coded by $Q_1, Q_0 = 10$. A correct string must start with 1, otherwise the automaton goes to the final state, $q_3$, which signals {\tt fail}. In state $q_1$, the automaton receives 1s until the first 0 is applied to the input, and the automaton passes to state $q_2$ in which it signals, through {\tt rec0}, that the string received is correct. If the state $q_2$ it is received in 1, then the string is cataloged as not belonging to those recognized by blocking the automaton in state $q_3$.

$\diamond$
\end{exemplu}

Designing a finite automaton using a description in System Verilog HDL involves writing the following files:
\begin{itemize}
\item {\tt defines.vh} in which the binary values that the state, input and output variables take are defined
\item {\tt topAutomaton.sv} in which the module that describes the automaton defining the state register, instantiating the combinational calculation module of the loop, {\tt loopCLC}, and the module {\tt outCLC} that describes the output circuit is defined
\item {\tt loopCLC.sv} which describes the state transition function
\item {\tt outCLC.sv} which describes the output transition function.
\end{itemize}


\begin{exemplu}\label{mealySV}
We now reconsider the previous example, which involves recognizing the language \( L = \{ a^n b^m \mid n, m > 0 \} \), and proceed with its implementation using a SystemVerilog description.

\normalfont{A key observation is that, regardless of the specific values of \( m \) and \( n \), the number of states in the finite automaton remains constant at four. However, while the number of states is fixed, the complexity of describing the automaton's behavior is proportional to this number. Specifically, the size of the truth table that defines the state transitions scales with the number of states but remains independent of the values of \( m \) and \( n \).}



\includeverilogcode[caption={Defines binary codes for input, output and state variables.  File: A10\_defines.vh  (SystemVerilog)}, 
                    label={lst:A10_defines.vh}, 
                    firstline=1, 
                    firstnumber=1, 
                    lastline=30]
                    {\verilogSourceDir/A10_defines.vh}

\end{exemplu}
Listing \ref{lst:A10_defines.vh} shows a code snippet that defines several preprocessor macros in SystemVerilog that are used to establish constants for a finite automaton. In the first section, input codes are defined. The macro \texttt{\`define a 1'b1} specifies that the symbol \textit{a} is represented by the 1-bit binary value \texttt{1'b1}, while \texttt{\`define b 1'b0} defines the symbol \textit{b} as the 1-bit binary value \texttt{1'b0}. These definitions are for the input signals to the automaton.

In the next portion of the code, output codes are established. The macro \texttt{\`define recA 2'b10} is used to indicate that the automaton outputs the 2-bit binary value \texttt{2'b10} when it receives input \textit{a}. Similarly, \texttt{\`define recB 2'b01} specifies that the output \texttt{2'b01} corresponds to the reception of input \textit{b}. Additionally, \texttt{\`define fail 2'b00} is defined to represent the output condition when the automaton fails to receive the correct string.

The final section of the code defines the internal states of the automaton. The initial state is set by the macro \texttt{\`define init 2'b10}, assigning the value \texttt{2'b10} to the starting state. The macro \texttt{\`define runA 2'b11} denotes a state where the automaton has received at least one \textit{b}. In contrast, \texttt{\`define runB 2'b01} is used to indicate that the automaton has received at least one \textit{a}. Finally, the macro \texttt{\`define error 2'b00} defines the error state, which shares the same binary encoding as the \texttt{fail} output.


\includeverilogcode[caption={Structural description of the immediate Mealy automaton
             designed to recognize the regular language $L = (a^nb^m | n,m >0)$.  File: A10\_regLang.sv  (SystemVerilog)}, 
                    label={lst:A10_regLang.sv}, 
                    firstline=1, 
                    firstnumber=1, 
                    lastline=30]
                    {\verilogSourceDir/A10_regLang.sv}

Listing \ref{lst:A10_regLang.sv} implements a Mealy automaton within a module named \texttt{regLang}. It begins by including a header file (\texttt{"defines.vh"}) that contains macro definitions for state encodings. The module itself declares an input signal (\texttt{in}), a two-bit output signal (\texttt{out}), a reset signal (\texttt{reset}), and a clock signal. Two internal two-bit registers, \texttt{state} and \texttt{nextState}, are defined to hold the current state and the computed next state of the automaton, respectively.

State updates occur in an \texttt{always\_ff} block that is triggered on the positive edge of the clock. Within this block, the current state is set to the initial state (defined by the macro \texttt{\`init}) when the reset signal is active. Otherwise, the state is updated to the value of \texttt{nextState}, which is determined by the automaton’s state-transition logic.

The next state and output logic are implemented in separate modules. The module instantiation \texttt{loopCLC lC( in, state, nextState )} is responsible for computing the next state based on the current state and the input signal. In parallel, the instantiation \texttt{outCLC oC( state, in, out )} computes the output, which depends on both the current state and the input signal. Because the output is a function of both the state and the input, this design characterizes a Mealy automaton, where outputs may change asynchronously with respect to the state update clock edge.

\includeverilogcode[caption={Combinational circuit used to compute the loop for the
                immediate Mealy automaton used to recognize the language $L = (a^nb^m | n,m >0)$.  File: A10\_loopCLC.sv  (SystemVerilog)}, 
                    label={lst:A10_loopCLC.sv}, 
                    firstline=1, 
                    firstnumber=1, 
                    lastline=30]
                    {\verilogSourceDir/A10_loopCLC.sv}

Listing \ref{lst:A10_loopCLC.sv} shows a SystemVerilog module with \texttt{loopCLC} that defines a combinational logic block that computes the next state of a finite state machine based on the current state and an input signal. The module accepts an input signal \texttt{in} and a two-bit signal \texttt{state} representing the current state, and it produces a two-bit output signal \texttt{nextState} that represents the computed next state. 

Within the module, an \texttt{always\_comb} block, labeled \texttt{loop\_clc}, implements the state transition logic using a \texttt{case} statement on the current state. When the state is \texttt{\`init}, the next state is determined by checking the input: if \texttt{in} equals the macro \texttt{\`a}, then \texttt{nextState} is assigned \texttt{\`runA}; otherwise, it is assigned \texttt{\`error}. In the \texttt{\`runA} state, if the input remains \texttt{\`a}, the automaton stays in the \texttt{\`runA} state, but if the input differs, it transitions to \texttt{\`runB}. When the state is \texttt{\`runB}, the next state remains \texttt{\`runB} if the input is equal to the macro \texttt{\`b}; if not, the automaton enters the \texttt{\`error} state. Once in the \texttt{\`error} state, the module is designed to remain in that state regardless of the input.

\includeverilogcode[caption={Combinational circuit used to compute the output for
             immediate Mealy  automaton which recognizes the language $L = (a^nb^m | n,m >0)$.   File: A10\_outCLC.sv  (SystemVerilog)}, 
                    label={lst:A10_outCLC.sv}, 
                    firstline=1, 
                    firstnumber=1, 
                    lastline=30]
                    {\verilogSourceDir/A10_outCLC.sv}


Listing \ref{lst:A10_outCLC.sv} shows a SystemVerilog module \texttt{outCLC} that implements the output logic for a finite state machine. It accepts a 2-bit input signal \texttt{state} and a 1-bit input \texttt{in}, and produces a 2-bit output signal \texttt{out}. This module uses an \texttt{always\_comb} block to define combinational logic that directly computes the output based on both the current state and the input value.

Within the combinational block, a \texttt{case} statement examines the value of \texttt{state} to determine the appropriate output. In the state defined by the macro \texttt{\`init}, the output is assigned conditionally: if the input \texttt{in} equals the macro \texttt{\`a}, then \texttt{out} is set to the value specified by \texttt{\`recA}; otherwise, it is set to \texttt{\`fail}. When the state is \texttt{\`runA}, if the input equals \texttt{\`a} the output remains \texttt{\`recA}, but if it does not, the output is updated to \texttt{\`recB}. For the state \texttt{\`runB}, the logic assigns \texttt{\`fail} to \texttt{out} when the input equals \texttt{\`a}, and assigns \texttt{\`recB} when it does not. Finally, if the state is \texttt{\`error}, the output is unconditionally set to \texttt{\`fail}.


\includeverilogcode[caption={Test module for the
             immediate Mealy  automaton which recognizes the language $L = (a^nb^m | n,m >0)$.   File: A10\_testRegLang.sv  (SystemVerilog)}, 
                    label={lst:A10_testRegLang.sv}, 
                    firstline=1, 
                    firstnumber=1, 
                    lastline=30]
                    {\verilogSourceDir/A10_testRegLang.sv}


Listing \ref{lst:A10_testRegLang.sv} shows a SystemVerilog module \texttt{testRegLang} that functions as a testbench for the finite state machine defined in the \texttt{regLang} module. In this testbench, several signals are declared, including a one-bit input \texttt{in}, a two-bit output \texttt{out}, a reset signal \texttt{reset}, and a clock signal \texttt{clock}. The testbench is responsible for generating stimulus and observing the behavior of the \texttt{regLang} module during simulation.

An initial block is dedicated to clock generation. The clock is initialized to 0, and a \texttt{forever} loop toggles the clock every time unit by inverting its value. This produces a periodic square wave that drives the synchronous components of the design. Another initial block manages the reset signal and input stimulus. Here, the reset is initially asserted, and the input is set to the value defined by the macro \texttt{\`a}. After a 2 time unit delay, the reset is de-asserted, and subsequent delays cause the input to change first to the value corresponding to \texttt{\`b} and later back to the value defined by \texttt{\`a}. The simulation is programmed to finish after 100 time units, ensuring that all intended transitions are exercised.

The \texttt{regLang} module, which contains the core finite state machine logic, is instantiated with the testbench signals connected to its ports. This instantiation, labeled as \texttt{dut} (device under test), allows the testbench to apply the generated inputs and observe the resulting outputs. Additionally, another initial block uses the \texttt{\$monitor} system task to continuously print the simulation time along with the current values of the clock, reset, input, and output signals, facilitating real-time observation of the system's behavior. Waveform dumping is also set up using \texttt{\$dumpfile} and \texttt{\$dumpvars}, which enable post-simulation analysis through a waveform viewer.

$\diamond$

\begin{verisum}{\em :

\hspace{-0.8cm} \rule[5mm]{16cm}{0.25mm} \vspace{-0.7cm}

\begin{description}
    \item[`define]: used to define global variable
    \item[`include] "fileName.vh" : used to include code already defined
    \item[`name]: used to assert a name already defined using {\tt `define}.
\end{description}
\rule[5mm]{16cm}{0.25mm} }
\end{verisum}


%****************************************************************
\subsection{{\bf {\em Control Automata (CROM)}}}\label{contrAut}
%****************************************************************
In addition to recognition and generation functions, finite state machines are used to control systems that evolve under sequences of received commands and return signals to the machine that report significant events for the evolution of the command sequence.

%*************************
\subsubsection{Structure}
%*************************
A control automaton is embedded within a system through three primary connections (see Figure \ref{contr_aut}):
\begin{itemize}
    \item The $p$-bit input {\tt operation[p-1:0]} selects the control
    sequence to be executed by the control automaton. This input is used 
    to partition the ROM into $2^{p}$ sections, each of equal size. 
    Within each section, a sequence of up to $2^{n}$ operations can be 
    stored for execution. Essentially, this input {\em specifies} "what to do."
    
    \item The $m$-bit command output {\tt command[m-1:0]} is used by 
    the control automaton to {\em generate} the necessary commands 
    for the controlled subsystem.
    
    \item The $n$-bit input {\tt flags[q-1:0]} allows the control 
    automaton to {\em receive} status information from the controlled 
    subsystem. This information consists of {\em independent bits} that 
    indicate "what is happening" within the subsystem in response to 
    the commands issued via {\tt command[m-1:0]}.
\end{itemize}


\placedrawing[h]{figures/04_04_control_automaton0.lp}{{\bf Control Automaton.} {\small The
functional definition of control automaton. Control means to issue
commands and to receive back signals (flags) characterizing the
effect of the command.}}{contr_aut}

The size and complexity of the control sequence necessitate the use of a ROM. The required ROM size is given by the following order of magnitude:

\[
S_{ROM}((m+n) \text{ outputs}, (n+p+q) \text{ inputs}) \in O((m+n) \times 2^{n+p+q}).
\]

To reduce the ROM size, we analyze real-world applications, which highlight two key observations:

\begin{enumerate}
    \item The automaton can {\em store} information about ``what to do'' within its state space. Specifically, each current state belongs to a path through the state space, {\bf originating} from an initial state determined by the code used to specify the operation.
    
    \item In most states, the automaton evaluates only a single bit from the {\tt flags[q-1:0]} input. In cases where additional bits must be considered, a small number of extra states in the flowchart typically resolve the issue.
\end{enumerate}


\placedrawing[h]{figures/04_05_control_automaton1.lp}{{\bf Optimized version of control
automata.} {\small The flags received from the controlled system
have independent meaning considered in distinct cycles. The flag
selected by the code TEST, T, decides from what half of ROM the
next state and output will be read.}}{red_ca}

Building upon these observations, the structure of the control automaton can be modified (see Figure \ref{red_ca}). Since the sequence is {\bf only} initialized using the code {\tt operation[n-1:0]}, this code is used solely to address the first command line in the ROM, within a single state where $MOD=1$. To implement this functionality, $n$ eMUXs must be added, along with a new ROM output to generate the signal $MOD$. This modification enables a more flexible utilization of the ``storage space'' within the ROM. Because the size of a control sequence can vary significantly from that of another, it is inefficient to allocate a fixed-size section of ROM for each sequence, as was done in the initial design. The revised structure shown in Figure \ref{red_ca} dynamically allocates only as much ROM space as necessary to store all command lines of each control sequence.

The second modification concerns the input {\tt flags[q-1:0]}. Since the bits associated with this input are tested in different states, $MUX_{q+2}$ selects the appropriate bit in each state using the $t$-bit field TEST. Consequently, the $q-1$ bits originally included in the {\tt flags[q-1:0]} input are removed from the ROM input, replacing them with only $t = \log_2 (q+2)$ output bits in the ROM. Instead of incorporating all $q$ bits, only a single bit, $T$, is connected to the ROM input. This new structure operates almost identically to the original but significantly reduces the required ROM size.

While working with this revised control automaton, another important observation emerges: {\em the majority of the generated sequence follows a linear structure.} As a result, the commands associated with these linear sequences can be stored in the ROM at consecutive addresses. That is, the next ROM address can be determined by incrementing the current address stored in register $R$. This leads to the structure illustrated in Figure \ref{crom}. The key innovation in this design is the inclusion of an increment circuit connected to the output of the state register, as well as a small combinational circuit that transcodes the bits $M_{1}, M_{0}, T$ into $S_{1}$ and $S_{0}$. There are four transition modes encoded by $M_{1}, M_{0}$:


\placedrawing[h]{figures/04_06_control_automaton2.lp}{{\bf The simplest Controller with
ROM (CROM).} {\small The Moore form of control automaton is
optimized using an incremented circuit (INC) to compute the most
frequent next address for ROM.}}{crom}

\begin{itemize}
\item {\tt inc}, encoded by {\tt MOD[1:0] = 00}: The next address 
for the ROM is obtained by incrementing the current address. 
The selection code must be {\tt S[1:0] = 00}.

\item {\tt jmp}, encoded by {\tt MOD[1:0] = 01}: The next address 
for the ROM is determined by the value stored in a specific field 
at the ROM output. The selection code must be {\tt S[1:0] = 01}.

\item {\tt cjmp}, encoded by {\tt MOD[1:0] = 10}: {\bf If} the 
selected flag, T, has a value of 1, {\bf then} the next address 
for the ROM is given by the value stored in a specific field at the 
ROM output; {\bf otherwise}, the next address is obtained by incrementing 
the current address. The selection code must be {\tt S[1:0] = 0T}.

\item {\tt init}, encoded by {\tt MOD[1:0] = 11}: The next address 
for the ROM is selected by $nMUX_{4}$ from the initialization input 
{\tt operation}. The selection code must be {\tt S[1:0] = 1-}.
\end{itemize}

This results in the following logic functions for the trans-coder TC:
\begin{quote}
{\tt S[1] = MOD[1]$\cdot$MOD[0]} \\ 
{\tt S[0] = MOD[1]$\cdot$T + MOD[0]}.
\end{quote}

The estimated size of the ROM is now given by:
\[
S_{ROM}((m+n+2+\log_2 (q+2)) \text{ outputs}, n \text{ inputs}) \in O((m+n) \times 2^{n+p+q})
\]
which indicates that the ROM size is reduced by a factor of 
$O(2^{p+q})$ compared to the initial unstructured version. Furthermore, 
all additional structures are simple circuits, meaning that the 
reduction is not only in size but, more importantly, in the complexity 
of the ROM’s content.

We refer to this optimized structure as CROM ({\bf C}ontroller with {\bf ROM}). 
CROM consists of both simple circuits (such as multiplexers, an incrementer, and registers) 
and more complex components (such as TC and ROM). The complexity of TC is negligible, 
as its size is minimal. The complexity of the ROM, however, is inherent in its content, 
which is defined as a patternless binary structure.

\subsubsection{Microcoding}

In order to define and verify the complex contents of the ROM, we will consider each location as a {\bf {\em microinstruction}} that is used to control the operation of effector systems. In this way the CROM will be defined by the {\bf {\em microprograms}} - the sequences of microinstructions - that it runs.
Thus, the output of ROM can be seen as a {\em microinstruction} defined
as follows: {\small
\begin{verbatim}
    <microinstruction>::= <<setLabel> <Mod> <Test> <Command>> ;
    <setLabel>::= lb(<number>);
    <number>::= 0 | 1 | ... | 9 | <number><number>;
    <mod>::=  jmp(<number>) | cjmp(<number>) |init | - ; // "-" means increment
    <test>::= <to be defined when use>;
    <command>::= <<to be defined when use> <endOfLine>>;
    <endOfLine>::= eol;
\end{verbatim}
}

\begin{exemplu}\label{dummyMP}
Let's customize the previous definition for a system that is commanded with 4 fields and returns to CROM two flags
 {\small
\begin{verbatim}
    <microinstruction>::= <<setLabel> <Mod> <Test> <Command>> ;
    <setLabel>::= lb(<number>);
    <number>::= 0 | 1 | ... | 9 | <number><number> | -;
    <mod>::=  jmp(<number>) | cjmp(<number>) |init | - ; // "-" means increment
    <test>::= flag1 | flag2 | - ;
    <command>::= <<field1> <field2> <field3> <field4> <endOfLine>>;
    <field1> ::= com11 | com12 | com13 | - ;
    <field2> ::= com21 | com22 | com23 | com24 | - ;
    <field3> ::= com31 | com32 | - ;
    <field4> :: <number>;
    <endOfLine>::= eol;
\end{verbatim}
}
Using the definition of {\tt microinstruction}, the content of the file {\tt theDefinition.v} defines, the following is a simple example microprogram:

{\em 
\begin{verilogcode}
/**********************************************************************
Function: Dummy microprogram
File name: theDefinition.sv
Description: a dummy microprogram used to test the code generator.
**********************************************************************/
    lb(1);  com13; com22; com31; val(24); eol;
            cjmp(3); flag2; com21; com32; eol;
            cjmp(1); flag1; eol;
    lb(7);  com31; val(42); eol;
            com12; com21; eol;
            cjmp(7); flag2; eol;
    lb(3);  jmp(3); eol;    // halt
\end{verilogcode}
}

$\diamond$
\end{exemplu}

{\bf A Very Important Comment!}

The CROM version of the control automaton's structure is characterized by two key processes:
\begin{itemize}
    \item An increase in structural complexity.
    \item A reduction in both the size and complexity of the 
    binary configuration ``stored'' in the ROM.
\end{itemize}

At this stage, both the size and complexity of the system increase without introducing any functional improvements. The sole effect is the reduction of the (algorithmic) complexity of the ROM's content.

This marks an important phase in the development of digital systems, where physical complexity begins to {\em compensate} for the ``symbolic'' complexity of the ROM's content. While both circuits and symbols are structured entities, there exists a fundamental difference between them. Physical structures can be defined through simple recursive relationships, whereas the symbolic content of the ROM is (almost) random and lacks a straightforward definition.

{\em We accept a moderate increase in the complexity of the physical 
structure, including its size, in order to create conditions that reduce the effort required to configure the complex symbolic content of the ROM.}

This represents the first major ``{\bf turning point}'' in the evolution of digital systems. Here, we encounter the initial indication of the greater complexity inherent in symbolic structures. By utilizing recursively defined components, physical structures can be kept relatively simple, whereas symbolic structures must account for the full complexity of the problems digital systems are designed to solve. 

The CROM structure is specifically designed so that the ROM content is {\em easy to design}, {\em easy to test}, and {\em easy to maintain}, despite its inherent complexity. This marks the first moment in our approach where the importance of the symbolic structure surpasses that of the physical structure in a digital system.


%***************************************************
\subsubsection{Binary code generator}\label{binGen}
%***************************************************


For defining the bit-content content of the ROM,
we have to design a Verilog description for a ``machine" which
takes a file containing lines of {\em microinstructions} and
translate it into the corresponding binary representation. Then,
in the module which describes the ROM of CROM we must include the following line: 
\begin{verilogcode}
  `include "codeGenerator.v" // generates ROM's content using 'theDefinition.v'
\end{verilogcode}

The code generating mechanism is a program which has as input a file describing the behavior of the automaton using a microprogram. Let us consider  Example \ref{dummyMP} where the file {\tt theDefinition} is used  as an input for {\tt codeGenerator.sv}  defined inListing \ref{lst:A10_codeGenerator.sv}.

\includeverilogcode[caption={Micro-program assembler..   File: A10\_codeGenerator.sv  (SystemVerilog)}, 
                    label={lst:A10_codeGenerator.sv}, 
                    firstline=1, 
                    firstnumber=1, 
                    lastline=130]
                    {\verilogSourceDir/A10_codeGenerator.sv}


Listing \ref{lst:A10_codeGenerator.sv} shows SystemVerilog code that generates the binary content of a ROM through a structured process that uses registers, tasks, and an initialization routine. Several registers are declared to represent various fields of the ROM entry. For instance, the 2-bit register \texttt{mode} indicates the transition mode, while \texttt{test} is a 1-bit register used for flag selection. Other registers—\texttt{field1} (2 bits), \texttt{field2} (3 bits), \texttt{field3} (2 bits), and \texttt{field4} (8 bits)—store different command parameters. The 5-bit registers \texttt{address} and \texttt{counter} keep track of the ROM addressing and entry count, respectively, and an array \texttt{labelTab} with 32 entries of 5 bits each is used to store addresses associated with specific label indices.

The code defines several tasks that facilitate the creation of ROM entries. The \texttt{eol} task, short for ``end of line'' finalizes an entry by combining all the defined fields into a single ROM word using a concatenation of \texttt{mode}, \texttt{address}, \texttt{test}, \texttt{field1}, \texttt{field2}, \texttt{field3}, and \texttt{field4}. After writing this composite value to the ROM at the location specified by \texttt{counter}, the task resets all fields to zero and increments the counter, preparing for the next entry.

A separate task, \texttt{lb}, is used for label management. It assigns the current value of \texttt{counter} to an index in the \texttt{labelTab} array, thereby enabling jump instructions to reference specific ROM addresses during program execution. In addition to these tasks, several command tasks (such as \texttt{com11}, \texttt{com12}, \texttt{com13} for modifying \texttt{field1}; \texttt{com21} through \texttt{com24} for modifying \texttt{field2}; \texttt{com31} and \texttt{com32} for modifying \texttt{field3}; and \texttt{val} for assigning a value to \texttt{field4}) are defined to set specific bits in the ROM entry corresponding to various instructions.

The code also includes tasks to define transition modes. For example, the \texttt{jmp} task sets \texttt{mode} to \texttt{2'b01} and loads \texttt{address} with a previously stored label from \texttt{labelTab}. The \texttt{cjmp} task functions similarly but sets \texttt{mode} to \texttt{2'b10} to support conditional jumps, while the \texttt{init} task, which sets \texttt{mode} to \texttt{2'b11}, indicates an initialization instruction.

Additionally, two tasks, \texttt{flag1} and \texttt{flag2}, are used for flag selection, assigning the values \texttt{1'b0} and \texttt{1'b1} to the \texttt{test} register, respectively.

The initialization block of the code sets all registers to zero and initializes the counter. It then includes an external file, \texttt{theDefinition.v}, twice. These include directives represent a two-pass assembly process: the first pass defines labels and records their corresponding addresses, while the second pass resolves these labels to generate the correct ROM content. The file {\tt theDefinition} is included twice because if a
label is used before it is defined, only during the second pass will the
memory {\tt labelTab} contain the correct value of the jump address.

\includeverilogcode[caption={Micro-program assembler..   File: A10\_simulator.sv  (SystemVerilog)}, 
                    label={lst:A10_simulator.sv}, 
                    firstline=1, 
                    firstnumber=1, 
                    lastline=130]
                    {\verilogSourceDir/A10_simulator.sv}


Listing \ref{lst:A10_simulator.sv} shows a SystemVerilog module, named \texttt{simulator}. It is a simulation framework that incorporates externally generated ROM content via the \texttt{`include} directive, which pulls in the file \texttt{codeGenerator.v}. Within the module, a ROM is declared as an array \texttt{mem} consisting of 32 elements, where each element is a 23-bit register. An integer variable \texttt{i} is also declared to serve as a loop index.

In the \texttt{initial} block, a for loop iterates from \texttt{i = 0} to \texttt{i < 16}, effectively processing the first 16 entries of the ROM. For each iteration, the \texttt{\$display} statement is executed, which prints the contents of the current ROM entry in a formatted string. Specifically, the output displays the ROM index, followed by several fields extracted from the corresponding 23-bit word. The fields are partitioned as follows: the \texttt{mode} field is obtained from bits 22 to 21, the \texttt{addr} field from bits 20 to 16, the \texttt{test} bit from bit 15, \texttt{com1} from bits 14 to 13, \texttt{com2} from bits 12 to 10, \texttt{com3} from bits 9 to 8, and finally the \texttt{val} field from bits 7 to 0, which is displayed as a decimal number.

This organization of bits within each ROM entry provides various control and command parameters, giving a structured way to interpret the ROM content generated by the code included from \texttt{codeGenerator.v}. Overall, the module functions as a simulator that outputs the ROM's binary data in a human-readable format, facilitating verification and debugging of the code generation process.

Running the simulation on the file {\tt theDefinition.x} results the content of the ROM:

\begin{verilogcode}
ROM[0] = mode:00 addr:00000 test:0 com1:11 com2:010 com3:01 val: 24
ROM[1] = mode:10 addr:00110 test:1 com1:00 com2:001 com3:10 val:  0
ROM[2] = mode:10 addr:00000 test:0 com1:00 com2:000 com3:00 val:  0
ROM[3] = mode:00 addr:00000 test:0 com1:00 com2:000 com3:01 val: 42
ROM[4] = mode:00 addr:00000 test:0 com1:10 com2:001 com3:00 val:  0
ROM[5] = mode:10 addr:00011 test:1 com1:00 com2:000 com3:00 val:  0
ROM[6] = mode:01 addr:00110 test:0 com1:00 com2:000 com3:00 val:  0
ROM[7] = mode:xx addr:xxxxx test:x com1:xx com2:xxx com3:xx val:  x
\end{verilogcode}

%%%%%%%%%%%%%%%%%%%%%%%%%%
\section{Simple Automata}
%%%%%%%%%%%%%%%%%%%%%%%%%%%
A simple automaton has the loop closed through a simple combinatorial circuit. We will provide two examples on our way to build a computing engine.

%*********************
\subsection{Counters}
%*********************

%***************************
\subsubsection{T Flip-Flop}
%***************************
The smallest and simplest automaton is a half-automaton with two states and a single input (see Figure \ref{t_ff}a). Given an automaton with only two internal states, denoted as \( Q \) and \( Q' \), and a single input bit \( T \), its behavior is constrained to a single logical operation:

\[
Q \Leftarrow T \;? \; Q' : Q
\]

\placedrawing[h]{figures/04_06_T_flip_flop.lp}{{\bf The T flip-flop.} {\small {\bf
a.} It represents the simplest form of an automaton, as it consists of a 1-bit state register (a D flip-flop), a two-input feedback loop (one input corresponding to the automaton's external input and the other maintaining the feedback loop), and an output directly derived from the state register. {\bf b.} The structural representation of the T flip-flop: the \( XOR_2 \) gate inverts the state when \( T = 1 \). {\bf c.} The logical symbol of the T flip-flop.}}{t_ff}

One fundamental property of the XOR function is its role as a controlled inverter. Consequently, the actual structure of the smallest and simplest automaton is illustrated in Figure \ref{t_ff}b.

The numerical function of the T flip-flop can be interpreted as a modulo-2 counter governed by the control input \( T \). When \( T = 1 \), the circuit toggles its state, effectively counting. Conversely, when \( T = 0 \), the circuit retains its previous state, exhibiting memory-like behavior.

%******************************
\subsubsection{Generic Counter}
%******************************
The generic counter is a basic half-automaton augmented with an increment circuit integrated into its feedback loop. At each active clock edge, the incremented value of the state register is loaded into the register. In practical applications, the counter's functionality is extended to support the following operations:

\placedrawing[h]{figures/04_07_counter.lp}{Counter.}{fig38}

\begin{itemize}
    \item {\tt op = 00 = nop: out <= out}
    \item {\tt op = 01 = reset: out <= 0}
    \item {\tt op = 10 = load: out <= in}
    \item {\tt op = 11 = count: out <= down ? out - 1 : out + 1}
    \item {\tt down}: 2's complement of {\tt out} to the increment circuit's input
\end{itemize}

The structure of the generic counter is represented in Figure \ref{fig38}. If the size of the register is of $n$ bits, then the circuit counts modulo $2^n$.

\begin{exemplu}
Four-bit presetable and reversible counter.
{\em 
\includeverilogcode[caption={Defines the command for the counter circuit.   File: A10\_counterDefines.vh  (SystemVerilog)}, 
                    label={lst:A10_counterDefines.vh}, 
                    firstline=1, 
                    firstnumber=1, 
                    lastline=30]
                    {\verilogSourceDir/A10_counterDefines.vh}
}

Listing \ref{lst:A10_counterDefines.vh} shows the commands for the counter automaton.
\end{exemplu}
 
\includeverilogcode[caption={A presetable up-down counter.   File: A10\_counter.sv  (SystemVerilog)}, 
                    label={lst:A10_counter.sv}, 
                    firstline=1, 
                    firstnumber=1, 
                    lastline=30]
                    {\verilogSourceDir/A10_counter.sv}

Listing \ref{lst:A10_counter.sv} shows SystemVerilog code that defines a module named \texttt{counter} that implements a simple synchronous counter. On line 1, the directive \texttt{`include "counterDefines.vh"} is used to incorporate a header file containing macro definitions. These macros, including \texttt{`reset}, \texttt{`load}, \texttt{`up}, and \texttt{`down}, specify the operation codes that determine the counter’s functionality.

The module \texttt{counter} has a 4-bit output, \texttt{out}, a 3-bit input, \texttt{op}, for selecting the operation, a 4-bit input, \texttt{in}, for loading a value, and a clock input, \texttt{clk}. The counter is designed to operate synchronously with the positive edge of the clock. Within an \texttt{always\_ff} block, which is triggered on each rising clock edge, a \texttt{case} statement evaluates the \texttt{op} input and updates the output accordingly.

When the \texttt{op} code corresponds to \texttt{`reset}, the counter is reset by setting \texttt{out} to 0. When the \texttt{op} code is \texttt{`load}, the value from the \texttt{in} input is loaded into \texttt{out}. For the \texttt{`up} operation, the counter increments its current value by 1, whereas for the \texttt{`down} operation, it decrements its current value by 1. In the default case, the counter maintains its current state by leaving \texttt{out} unchanged.


\includeverilogcode[caption={Test module for the counter.   File: A10\_testCounter.sv  (SystemVerilog)}, 
                    label={lst:A10_testCounter.sv}, 
                    firstline=1, 
                    firstnumber=1, 
                    lastline=30]
                    {\verilogSourceDir/A10_testCounter.sv}


Listing \ref{lst:A10_testCounter.sv} shows a SystemVerilog testbench module, \texttt{testCounter}, that is designed to verify the functionality of a counter module. On line 1, the directive \texttt{`include "counterDefines.vh"} incorporates an external header file that defines various operation macros (\texttt{`reset}, \texttt{`up}, \texttt{`load}, and \texttt{`down}).

Within the \texttt{testCounter} module, four signals are declared: a 4-bit output \texttt{out}, a 3-bit operation code \texttt{op}, a 4-bit input \texttt{in}, and a clock signal \texttt{clk}. The testbench generates a periodic clock by initializing \texttt{clk} to 0 and then continuously toggling its value every time unit using a \texttt{forever} loop. This clock generation mechanism ensures that the counter operates synchronously with a clock period of 2 time units.

Another initial block in the testbench establishes a sequence of test operations. Initially, the operation code \texttt{op} is set to \texttt{`reset}, and the input \texttt{in} is assigned the binary value \texttt{1010}. Following a delay of 2 time units, \texttt{op} is changed to \texttt{`up}, thereby instructing the counter to increment its current value. After an additional delay of 6 time units, the operation code is set to \texttt{`load}, causing the counter to load the value from \texttt{in}. Subsequently, after a further delay of 2 time units, the operation is again changed to \texttt{`up}, and then, following a delay of 10 time units, \texttt{op} is set to \texttt{`down} to prompt a decrement of the counter. Finally, after a delay of 16 time units, the simulation is terminated by invoking the \texttt{\$finish} system task.

At the conclusion of the testbench, the counter module is instantiated as the device under test (DUT) and connected to the declared signals \texttt{out}, \texttt{op}, \texttt{in}, and \texttt{clk}.

$\diamond$


%**********************************************************
\subsection{{\bf {\em Program Counter}}}\label{progCounter}
%**********************************************************

A specific application of the counter simple automaton is the system used to compute the next address in the process of running a program (see Figure \ref{fig39}). The set of functions specified on the input {\tt next} are:
\begin{description}
    \item[nop]: for {\tt halt} instruction the program counter stays unchanged \\ 
        {\tt  address <= address} 
    \item[rst]: for system reset the address is initialized to zero \\
        {\tt  address <= 0} 
    \item[inc]: for program counter increment \\
        {\tt  address <= address + 1} 
    \item[rjmp]: for relative jump \\
        {\tt address <= address + relValue} 
    \item[crjmp]: for conditioned branch \\
        {\tt  address <= (value = 0) ? address + relValue : address + 1} 
    \item[ajmp]: for absolute jump \\
        {\tt  address <= absValue} 
\end{description}

\placedrawing[H]{figures/04_08_program_counter.lp}{Program Counter.}{fig39}

\includeverilogcode[caption={Program counter.   File: A10\_programCounter.sv  (SystemVerilog)}, 
                    label={lst:A10_programCounter.sv}, 
                    firstline=1, 
                    firstnumber=1, 
                    lastline=30]
                    {\verilogSourceDir/A10_programCounter.sv}

Listing \ref{lst:A10_programCounter.sv} shows SystemVerilog code for a module named \texttt{programCounter} that functions as a program counter with a parameterized width. The module outputs a 32-bit program address (\texttt{progAddr}) and receives several inputs: a 3-bit selection signal (\texttt{sel}), a 32-bit relative value (\texttt{relValue}), a 32-bit absolute value (\texttt{absValue}), as well as \texttt{reset} and \texttt{clock} signals. Internally, two 32-bit registers, \texttt{pc} and \texttt{nextPC}, are declared; \texttt{pc} holds the current program counter value, while \texttt{nextPC} is computed based on the control inputs.

The sequential behavior of the counter is described within an \texttt{always\_ff} block that is triggered on the positive edge of the clock. Within this block, if the \texttt{reset} signal is asserted, the program counter \texttt{pc} is reset to zero; otherwise, \texttt{pc} is updated with the value of \texttt{nextPC}. This ensures that the program counter is synchronously updated and can be reliably reset when needed.

The combinational logic for computing the next program counter value is implemented in an \texttt{always\_comb} block using a \texttt{case} statement based on the \texttt{sel} input. When \texttt{sel} is \texttt{3'b000}, \texttt{nextPC} is assigned the current value of \texttt{pc}, effectively holding the counter. If \texttt{sel} is \texttt{3'b001}, \texttt{nextPC} is computed as \(\texttt{pc} + \texttt{relValue}\). For a selection of \texttt{3'b010}, the value of \texttt{nextPC} depends on the absolute value: if \texttt{absValue} is zero, \(\texttt{nextPC} = \texttt{pc} + \texttt{relValue}\); otherwise, \(\texttt{nextPC} = \texttt{pc} + 1\). Similarly, when \texttt{sel} is \texttt{3'b011}, \(\texttt{nextPC}\) is assigned \(\texttt{pc} + 1\) if \texttt{absValue} is zero, or \(\texttt{pc} + \texttt{relValue}\) otherwise. If \texttt{sel} is \texttt{3'b100}, the counter is loaded with the absolute value, so \(\texttt{nextPC} = \texttt{absValue}\). For any other value of \texttt{sel} (the default case), \(\texttt{nextPC} = \texttt{pc} + 1\).

Finally, the current value of the program counter \texttt{pc} is assigned to the output \texttt{progAddr}. An alternative assignment (commented out in the code) would directly assign \texttt{nextPC} to \texttt{progAddr} for a non-pipelined implementation.

%***********************************************************************************
\subsection{{\bf {\em Registers with Arithmetic \& Logic Unit (RALU)}}}\label{rALU}
%***********************************************************************************

If we connect in a loop a register file with a arithmetic logic unit we obtain the simplest form of an executive core of a processor (see Figure \ref{ralu}). The external connections of a Register with Arithmetic \& Logic Unit (RALU) type module are:
\begin{description}
    \item[{\tt func[3:0]}]: selects one of the maximum 8 functions performed by ALU
    \item[{\tt carryIn}]: carry input for arithmetic functions (is borrow for subtract)
    \item[{\tt carryOut}]: carry output for arithmetic function (is borrow for subtract)
    \item[{\tt in[31:0]}]: data input
    \item[{\tt load}]: selects data input as left operand for ALU
    \item[{\tt leftAddr[4:0]}]: selects left output or left operand for ALU when {\tt load = 0}
    \item[{\tt rightAddr[4:0]}]: selects right output or right operand for ALU
    \item[{\tt clock}]: is the clock signal active on the positive edge
    \item[{\tt we}]: write enable for the register file
    \item[{\tt destAddr[4:0]}]: destination address for the value {\tt out} provided by ALU
    \item[{\tt leftOut[31:0]}]: left output of RALU
    \item[{\tt rightOut[31:0]}]: right output of RALU
\end{description}
According to the connection list, RALU performs no more than 8 functions, on 32-bit words stored in a 32-word register file.

\placedrawing[h]{figures/04_09_ralu.lp}{32-bit RALU.}{ralu}


\includeverilogcode[caption={Defining the functions for a RALU.   File: A10\_defineRALU.vh  (SystemVerilog)}, 
                    label={lst:A10_defineRALU.vh}, 
                    firstline=1, 
                    firstnumber=1, 
                    lastline=30]
                    {\verilogSourceDir/A10_defineRALU.vh}



\includeverilogcode[caption={RALU.   File: A10\_RALU.sv  (SystemVerilog)}, 
                    label={lst:A10_RALU.sv}, 
                    firstline=1, 
                    firstnumber=1, 
                    lastline=30]
                    {\verilogSourceDir/A10_RALU.sv}


Listing \ref{lst:A10_defineRALU.vh} shows SystemVerilog code for a module named \texttt{RALU}. It implements a Register Arithmetic Logic Unit by integrating a register file and an arithmetic logic unit (ALU). The module is designed for 32-bit wide data processing and features three primary outputs: two 32-bit signals, \texttt{leftOut} and \texttt{righttOut}, and a one-bit flag output, \texttt{crOut}. The inputs include three 5-bit addresses (\texttt{leftAddr}, \texttt{rightAddr}, and \texttt{destAddr}) used to access specific registers, a 32-bit data input \texttt{in}, a control signal \texttt{load} that selects between direct input and register file output, a write enable signal \texttt{we}, a 4-bit function code \texttt{func} that determines the ALU operation, and a clock signal.

Internally, the module declares a 32-bit logic variable \texttt{out}, which serves as the output from the ALU and as the data input to the register file. The module instantiates a register file,  \texttt{rf}, which manages data storage and retrieval. For module {\tt regFile} see Section \ref{registerFile}. The outputs of the register file are connected to \texttt{leftOut} and \texttt{righttOut}, while the ALU’s result, \texttt{out}, is routed back into the register file for potential writing. The addresses for reading from the register file are supplied by \texttt{leftAddr} and \texttt{rightAddr}, and the destination register is specified by \texttt{destAddr}. Additionally, the write enable signal \texttt{we} and the clock signal are forwarded to the register file to control synchronous data updates.

Simultaneously, the module instantiates an ALU, denoted as \texttt{alu}, which performs arithmetic or logical operations based on the provided function code \texttt{func}. The left input to the ALU is determined by a conditional expression: if the \texttt{load} signal is asserted, the left input is taken directly from the external 32-bit input \texttt{in}; otherwise, it is taken from \texttt{leftOut}, the output of the register file. The right input to the ALU is directly connected to \texttt{righttOut}. The ALU computes its result and outputs it to the internal variable \texttt{out}, while also generating the carry or flag output \texttt{crOut}.

\includeverilogcode[caption={ALU.   File: A10\_alu.sv  (SystemVerilog)}, 
                    label={lst:A10_alu.sv}, 
                    firstline=1, 
                    firstnumber=1, 
                    lastline=30]
                    {\verilogSourceDir/A10_alu.sv}

Listing \ref{lst:A10_alu.sv} shows SystemVerilog code that defines an Arithmetic Logic Unit (ALU) module that performs a variety of arithmetic and logical operations on 32-bit operands. The module begins by including an external file, \texttt{"define.vh"}, which contains macro definitions for the operation codes used within the design. The ALU accepts a 4-bit function code, \texttt{func}, along with two 32-bit inputs, \texttt{left} and \texttt{right}. It produces a 32-bit output, \texttt{out}, and a 1-bit carry or flag output, \texttt{crOut}.

Two internal 33-bit signals, \texttt{sum} and \texttt{dif}, are declared to hold the results of addition and subtraction, respectively. The extra bit in these signals is used to capture any carry-out that may occur during arithmetic operations. Specifically, \texttt{sum} is assigned the result of \texttt{left + right}, while \texttt{dif} is assigned the result of \texttt{left - right}.

The core functionality of the ALU is implemented within an \texttt{always\_comb} block that employs a \texttt{case} statement to select the appropriate operation based on the value of \texttt{func}. For instance, when the function code corresponds to \texttt{`bigv}, the module concatenates a zero bit with the lower 16 bits of \texttt{left} and \texttt{right}, effectively merging these halves into a 32-bit output while setting the carry flag to zero. In the cases of \texttt{`add} and \texttt{`sub}, the concatenation of the carry-out and the 32-bit result from the addition (\texttt{sum}) or subtraction (\texttt{dif}) is assigned to the outputs.

Additional operations are provided for the conditional output of the carry bit. When the function code is \texttt{`addcr} or \texttt{`subcr}, the ALU outputs only the carry bit from the corresponding arithmetic operation, with the remaining 32 bits set to zero. The ALU also supports logical and shifting operations. The \texttt{`lsh} and \texttt{`ash} operations perform logical and arithmetic shifts on the \texttt{left} operand, with the shifted-out bit being routed to the carry output. Other operations, such as \texttt{`move}, \texttt{`mult}, \texttt{`bwand}, \texttt{`bwor}, and \texttt{`bwxor}, respectively, transfer the \texttt{left} operand, multiply the operands, and execute bitwise AND, OR, and XOR functions, with the carry output remaining zero in these cases.

In the default case, the outputs are assigned a 33-bit value corresponding to zero minus one, effectively setting all bits of \texttt{out} and \texttt{crOut} to zero, except for the bitwise representation of negative one. 


Let us use the RALU unit we just defined to illustrate elementary computational operations.

\begin{exemplu}\label{ex3.1}
The following sequence of operations: load two numbers (12 and 5) in the register file, add them, divide the result by 4, add with 23, and send out the result on the right output.


\begin{table}[H]
  \begin{center}
    \caption{Sequence of commands applied to RALU in Example \ref{ex3.1}}
    \label{raluSeq}
    \vspace{2mm}
    %{\small
    {\tt
    \begin{tabular}{|c|c|c|c|c|c|c||c|}
    \hline
    func &  in & load & leftAddr & rightAddr & we & destAddr & COMMENT          \\ \hline \hline
    1111 &  xx &  0   &  00000   &  00000    & 1  &  00000   & R0 <= 0          \\ \hline
    0010 &  12 &  1   &  xxxxx   &  00000    & 1  &  00001   & R1 <= 12+R0      \\ \hline
    0010 &   5 &  1   &  xxxxx   &  00000    & 1  &  00010   & R2 <= 5+R0       \\ \hline
    0010 &  xx &  0   &  00001   &  00010    & 1  &  00011   & R3 <= R1+R2      \\ \hline
    1001 &  xx &  0   &  00011   &  xxxxx    & 1  &  00011   & R3 <= R3/2       \\ \hline
    1001 &  xx &  0   &  00011   &  xxxxx    & 1  &  00011   & R3 <= R3/2       \\ \hline
    0010 &  23 &  1   &  xxxxx   &  00011    & 1  &  00011   & R3 <= R3+23      \\ \hline
    xxxx &  xx &  x   &  xxxxx   &  00011    & 0  &  xxxxx   & rightOut = R3    \\ \hline
 \end{tabular}
 }
 %}
  \end{center}
\end{table}

$\diamond$
\end{exemplu}


A more flexible way to test the correctness of the design is to make a simulation.

\includeverilogcode[caption={Testing RALU.   File: A10\_testRALU.sv  (SystemVerilog)}, 
                    label={lst:A10_testRALU.sv}, 
                    firstline=1, 
                    firstnumber=1, 
                    lastline=60]
                    {\verilogSourceDir/A10_testRALU.sv}


The result of simulation is printed by the monitor defined in the previous test module:

{\footnotesize
\begin{verbatim}
t=0  clock=0 func=1010 destAddr=0 leftAddr=6 rightAddr=7 load=1 in=8  RF=[x, x, x, x]
t=1  clock=1 func=1010 destAddr=0 leftAddr=6 rightAddr=7 load=1 in=8  RF=[8, x, x, x]
t=2  clock=0 func=1010 destAddr=1 leftAddr=2 rightAddr=2 load=1 in=9  RF=[8, x, x, x]
t=3  clock=1 func=1010 destAddr=1 leftAddr=2 rightAddr=2 load=1 in=9  RF=[8, 9, x, x]
t=4  clock=0 func=1010 destAddr=2 leftAddr=6 rightAddr=7 load=1 in=10 RF=[8, 9, x, x]
t=5  clock=1 func=1010 destAddr=2 leftAddr=6 rightAddr=7 load=1 in=10 RF=[8, 9, 10, x]
t=6  clock=0 func=1010 destAddr=3 leftAddr=2 rightAddr=2 load=1 in=11 RF=[8, 9, 10, x]
t=7  clock=1 func=1010 destAddr=3 leftAddr=2 rightAddr=2 load=1 in=11 RF=[8, 9, 10, 11]
t=8  clock=0 func=0010 destAddr=3 leftAddr=2 rightAddr=1 load=0 in=7  RF=[8, 9, 10, 11]
t=9  clock=1 func=0010 destAddr=3 leftAddr=2 rightAddr=1 load=0 in=7  RF=[8, 9, 10, 19]
t=10 clock=0 func=0011 destAddr=3 leftAddr=2 rightAddr=1 load=0 in=7  RF=[8, 9, 10, 19]
t=11 clock=1 func=0011 destAddr=3 leftAddr=2 rightAddr=1 load=0 in=7  RF=[8, 9, 10, 1]
t=12 clock=0 func=1101 destAddr=3 leftAddr=2 rightAddr=1 load=0 in=7  RF=[8, 9, 10, 1]
t=13 clock=1 func=1101 destAddr=3 leftAddr=2 rightAddr=1 load=0 in=7  RF=[8, 9, 10, 8]
\end{verbatim}
}

%****************************
\subsection{Stream Memories}
%****************************
One of the most significant applications of counters is to facilitate addressing within RAM structures, thereby enabling the systematic organization and retrieval of data in the form of strings. At both the conceptual level, within software, and the structural level, within hardware, two distinct architectures for storing data strings have been devised. One architecture permits both writing and reading operations at the same end of the string, whereas the other supports writing at one end and reading at the opposite end.

The first architecture corresponds to Last-In-First-Out (LIFO) memory, commonly referred to as a {\bf {\em stack}}, while the second corresponds to First-In-First-Out (FIFO) memory, also known as a {\bf {\em queue}}.





%**************************************************
\subsubsection{LIFO (Stack) Memory}\label{lb:lifo}
%**************************************************
\paragraph{Definition}  The definition of LIFO, or stack memory, considers a sequence of symbols, denoted by \texttt{S = <s0, s1, s2, ...>}, which evolves according to a series of operations.  For example, the following sequence illustrates the typical behavior: 
\begin{verbatim}
    push a --> S = <a, s0, s1, s2, ...>
    push b --> S = <b, a, s0, s1, s2, ...>
    pop    --> S = <a, s0, s1, s2, ...>
    pop    --> S = <s0, s1, s2, ...>
    pop    --> S = <s1, s2, ...>
\end{verbatim}

In practical applications, especially in expression evaluation, additional functionalities are required for a LIFO memory. A minimally extended set of operations for the LIFO, represented as \texttt{S = <s0, s1, s2, ...>}, comprises the following commands:
\begin{itemize}
\item {\tt nop}: no operation; the contents of {\tt S} in unchanged
\item {\tt write a}: write {\tt a} in {\tt TOS} \\{\tt <s0, s1, s2, ...> --> <a, s1, s2, ...>}\\used for unary operations; for example: \\{\tt a = s0 + 1}\\
the TOS is incremented and write back in TOS ({\tt pop, push = write})
\item {\tt pop}: \\{\tt <s0, s1, s2, ...> --> <s1, s2, ...>}
\item {\tt popwr a}: pop \& write \\{\tt <s0, s1, s2, ...> --> <a, s2, ...>}\\used for binary operations; for example: \\{\tt a = s0 + s1}\\
the first two positions in LIFO are popped, added and the result is pushed back into the LIFO memory ({\tt pop, pop, push = pop, write})
\item {\tt push a}: \\{\tt <s0, s1, s2, ...> --> <a, s0, s1, s2, ...>}
\end{itemize}

\placedrawing[H]{figures/so2_10.lp}{{\bf Stack memory}. {\small {\bf a.} The structure. {\bf b.} Logic symbol. }}{stack}


\paragraph{Structure} of a  LIFO memory is defined in Listing \ref{lst:A10_lifo.sv}.

\includeverilogcode[caption={LIFO.   File: A10\_lifo.sv  (SystemVerilog)}, 
                    label={lst:A10_lifo.sv}, 
                    firstline=1, 
                    firstnumber=1, 
                    lastline=60]
                    {\verilogSourceDir/A10_lifo.sv}

Listing \ref{lst:A10_lifo.sv} shows a SystemVerilog module that implements a LIFO memory structure, commonly known as a stack. The module, named \texttt{lifo}, features two 32-bit outputs, \texttt{stack0} and \texttt{stack1}, which are derived from an underlying register file. It accepts a 32-bit data input (\texttt{in}), a 3-bit command input (\texttt{com}), and control signals for reset and clock.

Within the module, command codes are defined by convention (as indicated in the comments) to control stack operations. These commands include \texttt{nop (3'b000)} for no operation, \texttt{write (3'b001)} to enable writing, \texttt{pop (3'b010)} to decrement the stack pointer, \texttt{popwr (3'b011)} to perform a combined decrement and write, and \texttt{push (3'b101)} to increment the pointer while writing. 

A 5-bit signal, \texttt{leftAddr}, serves as the main pointer to the current top of the stack, while \texttt{nextAddr} represents the computed subsequent pointer value. The assignment to \texttt{nextAddr} is determined by the command input: if \texttt{com[2]} is asserted, \texttt{nextAddr} is set to \texttt{leftAddr + 1}; if \texttt{com[2]} is not asserted but \texttt{com[1]} is, \texttt{nextAddr} is set to \texttt{leftAddr - 1}; otherwise, \texttt{leftAddr} remains unchanged. An \texttt{always\_ff} block, triggered on the positive edge of the clock, updates \texttt{leftAddr} accordingly—resetting it to zero when the \texttt{reset} signal is active, or updating it with the value of \texttt{nextAddr} otherwise.

The module then instantiates a register file (named \texttt{rf}) that manages the actual data storage. In this instantiation, \texttt{leftAddr} is used as the address for the left output port, while \texttt{leftAddr - 1} is used for the right output port. The destination address for data writing is provided by \texttt{nextAddr}, and the write enable signal is extracted from \texttt{com[0]}. The input data (\texttt{in}) is routed into the register file, and the entire operation is synchronized with the clock signal.

It is important to address the question of why an explicit signal to indicate the maximum storage capacity of the stack is not provided. In practice, the depth of the stack is typically designed to be sufficiently large such that overflow conditions are highly unlikely to occur. Nonetheless, mechanisms to signal when the storage capacity is exceeded do exist and can be implemented as needed.

%**************************************************
\subsubsection{FIFO (Queue) Memory}\label{lb:fifo}
%**************************************************

\paragraph{Definition} of the FIFO memory is shown in Figure \ref{fifo1}. It has the following connections (see Figure \ref{fifo1}a):
%\begin{description}[parsep=0em, itemsep=0.25em]
\begin{description}
    \item[dataIn]: the data entry of $n$ bits to which the word that will be recorded in the rightmost unoccupied location is applied.
    \item[dataOut]: the $n$ bit data output from which the oldest unextracted record is extracted
    \item[full]: it signals the fact that the memory is full and will not accept a new write before executing at least one read
    
    \item[write]: if activated when {\tt full = 0}, then {\tt dataIn} is queued
    \item[empty]: it signals the fact that the memory is empty and will not accept a new read before executing at least one write
    \item[read]: if activated when {\tt empty = 0}, then {\tt dataOut} is extracted from the  queue
\end{description}

In Figures from \ref{fifo1}b to \ref{fifo1}e is exemplified the way FIFO works, as follows:
\begin{description}
    \item[{\tt write(A)}]: in the empty FIFO A is stored in the rightmost position (see Figure \ref{fifo1}b)
    \item[{\tt write(D)}]: B is stored in the rightmost free position, immediately after A (see Figure \ref{fifo1}c)
    \item[{\tt write(F)}]: F is stored in the rightmost free position, immediately after B (see Figure \ref{fifo1}cd)
    \item[{\tt read}]: A is extracted from FIFO and the queue "advances" one step (see Figure \ref{fifo1}e)
\end{description}

\placedrawing[H]{figures/06_10_fifo.lp}{{\bf a.} Block schematic of FIFO. {\bf b.} {\tt write A} operation. {\bf c.} {\tt write D.} {\bf d.} {\tt write F}. {\bf e.} {\tt read}.}{fifo1}

There are two main types of FIFO:
\begin{itemize}
    \item synchronous FIFO: the sending system and the receiving system run with the same clock signal 
    %(for those who are interested can consult Appendix \ref{fifoMem})
    
    \item asynchronous FIFO: the sending system and the receiving system run with different clock signals
\end{itemize}
Asynchronous FIFOs are used to solve the problem of crossing data between systems running in different clock domains. The module BUS Bridge (see Figure \ref{compSys}) contains an asynchronous FIFO.

\paragraph{Structure} of a synchronous FIFO memory is shown in Figure \ref{fifo}, where:

\begin{itemize}
    \item {\bf RAM} is a $2^{n}$ $m$-bit words two-port asynchronous random access memory, one
    port for write to the address {\tt w\_addr} and another for
    read form the address {\tt r\_addr}
    \item {\bf write counter} is an $(n+1)$-bit resetable counter incremented
    each time a write is executed; its output is {\tt w\_addr[n:0]}, initially
    it is reset
    \item {\bf read counter} is an $(n+1)$-bit resetable counter incremented
    each time a read is executed; its output is {\tt r\_addr[n:0]}, initially it is reset
    \item {\bf eq} is a comparator activating its output when
    the least significant $n$ bits of the two counters are
    identical.
\end{itemize}

\placedrawing[H]{figures/draw0_28.lp}{{\bf FIFO memory.} {\small Two
pointers, evolving in the same direction, and a two-port RAM
implement a LIFO (queue) memory. The limit flags are computed
combinational from the addresses used to write and to read the
memory.}}{fifo}




A FIFO memory operates as a circular buffer, where data is addressed by two pointers: the write pointer, {\tt w\_addr[n-1:0]}, and the read pointer, {\tt r\_addr[n-1:0]}. These pointers progress in the same direction, ensuring sequential access to stored data. The memory is considered \textit{full} when, following a write operation, the write pointer becomes equal to the read pointer. In this case, the {\tt full} signal is asserted (\texttt{full = 1}). Conversely, the memory is deemed \textit{empty} when, after a read operation, the read pointer becomes equal to the write pointer, in which case the {\tt empty} signal is asserted (\texttt{empty = 1}).

To distinguish between the full and empty states when both pointers have the same value, an additional bit, the $(n+1)$-th bit, is utilized in each counter. If {\tt w\_addr[n]} and {\tt r\_addr[n]} differ, then {\tt w\_addr[n-1:0]} being equal to {\tt r\_addr[n-1:0]} indicates a \textit{full} condition. Otherwise, when {\tt w\_addr[n]} and {\tt r\_addr[n]} are the same, this equality signifies an \textit{empty} condition.

The comparison between the two address pointers is performed by a purely combinational circuit, which introduces potential timing hazards affecting the {\tt full} and {\tt empty} outputs. As a result, these signals must be used with caution when designing systems that incorporate this FIFO memory. However, this issue is mitigated because the entire system operates within the same clock domain, ensuring that both ends of the FIFO and the rest of the system are synchronized to the same {\tt clock} signal. A FIFO implementation operating under this condition is referred to as a \textit{synchronous FIFO}.

%%%%%%%%%%%%%%%%%%%%
\section{Problems}
%%%%%%%%%%%%%%%%%%%%%
\subsection{Finite Complex Automata}

\begin{problem}
Design the finite automaton which recognise the binary strings $a^{n}b^2c^{m}$, for $n \geq 1$ and $m \geq 1$. The input and the output signal are specified in Listing \ref{lst:A10_defines1.vh} (where the state assignment must be added):
{\em
\includeverilogcode[caption={Binary codes for input, output and state variables.   File: A10\_defines1.vh  (SystemVerilog)}, 
                    label={lst:A10_defines1.vh}, 
                    firstline=1, 
                    firstnumber=1, 
                    lastline=60]
                    {\verilogSourceDir/A10_defines1.vh}
}

$\diamond$
\end{problem}

\begin{problem}
Design the interrupt automaton which manage the interrupt process in a processor. The synchronous interaction of a processor with an external event, in its simplest form, is achieved by accepting an interrupt signal, {\tt int}, to which it responds with {\tt inta} ("interrupt acknowledgment"). Acceptance of the interruption is conditioned by the state {\tt enableState} state in which the interrupt management automaton can be. At {\tt reset}, the automaton goes into the state {\tt disableState}. The {\tt eint} instruction puts the automaton in the {\tt enableState} state, a state in which {\tt inta} can be issued if {\tt int} is received. After issuing {\tt inta}, the automaton generates the signal {\tt flush} which blocks the update of the processor state for a number of cycles, number that depends on the depth of the execution pipe (3, in the considered case). At the end of {\tt flush} signal the automaton switches in the  {\tt disableState}. Instruction {\tt dint} sends the automaton in the {\tt disableState} state.

{\em
\includeverilogcode[caption={Binary codes for input, output and state variables.   File: A10\_defines2.vh  (SystemVerilog)}, 
                    label={lst:A10_defines2.vh}, 
                    firstline=1, 
                    firstnumber=1, 
                    lastline=60]
                    {\verilogSourceDir/A10_defines2.vh}
}     

\begin{enumerate}
    \item Draw the graph of the automaton whose external connections are partially defined as follows:
Drawing the graph will allow completing the previous definitions file.
    \item Provide the System Verilog description of the automaton
    \item Use the Vivado environment to draw the schematic of the resulting circuit
    \item Simulate the automaton to test its behavior.
\end{enumerate}

$\diamond$
\end{problem}

\begin{problem}
Let be the automaton defined by the graph represented in Figure \ref{graph}.

\placedrawing[h]{figures/04_10_graph.lp}{}{graph}

{\em
\includeverilogcode[caption={Binary codes for input, output and state variables.   File: A10\_defines3.vh  (SystemVerilog)}, 
                    label={lst:A10_defines3.vh}, 
                    firstline=1, 
                    firstnumber=1, 
                    lastline=60]
                    {\verilogSourceDir/A10_defines3.vh}
}    
Implement the automaton in two versions assigning for the 5 states the binary codes in different versions. Compare the size of the two solutions.

$\diamond$
\end{problem}

\subsection{Simple Functional Automata}

\begin{problem}
Consider an 8-bit counter whose counting modes are selected by three bits as follows:
\begin{enumerate}
    \item no operation
    \item increment modulo $2^8$ by one unit at each clock cycles
    \item increment modulo $2^8$ by 3 at each clock cycles
    \item increment modulo $2^8$ by 5 at each clock cycles
    \item increment modulo $2^8$ by 7 at each clock cycles
    \item increment modulo $2^8$ by 11 at each clock cycles
    \item increment modulo $2^8$ by 13 at each clock cycles
    \item increment modulo $2^8$ by 17 at each clock cycles
\end{enumerate}
Required:
\begin{enumerate}
    \item Draw the block diagram of the solution you propose using known circuits
    \item Write the structural System Verilog description of the circuit
    \item Simulate the design to prove the functionality
\end{enumerate}

$\diamond$
\end{problem}

\begin{problem}
Simulate the circuit described in Section \ref{progCounter}.

$\diamond$
\end{problem}

\begin{problem}
Consider the program counter described in Section \ref{progCounter}.

1. Draw the logic diagram of the circuit described in the {\tt programCounter.sv} module using a register, and studied combinational circuits.

2. If $t_{pReg} = \#30$, $t_{SU} = \#5$, $t_{selectionMUX} = \#15$, $t_{selectedMUX} = \#5$, $t_{inc} = \#100$ , $t_{add} = \#110$, then what is the minimum period of the clock signal.

$\diamond$
\end{problem}

\begin{problem}
Simulate the circuit described in Section \ref{rALU} running the sequence outlined in Table \ref{raluSeq}.

$\diamond$
\end{problem}

\begin{problem}
Upgrade the RALU described in Section \ref{rALU} with ALU upgraded in Problem \ref{upgradeALU}.

$\diamond$
\end{problem}


\begin{problem}
Add to the RALU described in Section \ref{rALU} the possibility that the right operand can come from a second selected input with a distinct associated signal. Instead of {\tt in[31:0]} and {\tt load}, the design will contain {\tt inLeft[31:0]}, {\tt loadLeft}, {\tt inRight[31:0]} and {\tt loadRight}.

$\diamond$
\end{problem}


\begin{problem}
Consider the RALU described in Section \ref{rALU}, where: $t_{accessRegFile} = \#50$, $t_{rfSU} = \#5$, $t_{selectionMUX} = \#15$, $t_{selectedMUX} = \#5$, $t_{pALU} = \#150$. What is the minimum period of the clock signal.

$\diamond$
\end{problem}


\begin{problem}
Consider the RALU described in Section \ref{rALU} and fill-up the first 7 columns in Table \ref{aluSeq} according to operations specified in the last column.

\begin{table}[H]
  \begin{center}
    \caption{Sequence of commands applied to RALU described in Section \ref{rALU}}
    \label{aluSeq}
    \vspace{2mm}
    %{\small
    {\tt
    \begin{tabular}{|c|c|c|c|c|c|c||c|}
    \hline
    func &  in & load & leftAddr & rightAddr & we & destAddr & COMMENT          \\ \hline \hline
     &   &     &     &      &   &     & R2 <= 5         \\ \hline
     &   &     &     &      &   &     & R5 <= 23        \\ \hline
     &   &     &     &      &   &     & R12 <= R2 + R3 	\\ \hline
     &   &     &     &      &   &     & R14 <= R12      \\ \hline
     &   &     &     &      &   &     & R1 <= R14       \\ \hline
     &   &     &     &      &   &     & R2 <= R5 + R13  \\ \hline
     &   &     &     &      &   &     & R7 <= R2/2      \\ \hline
     &   &     &     &      &   &     & R8 <= R7 +12    \\ \hline
 \end{tabular}
 }
 %}
  \end{center}
\end{table}


$\diamond$
\end{problem}

\begin{problem}
Simulate the LIFO memory described in Listing \ref{lst:A10_lifo.sv}.

$\diamond$
\end{problem}

\begin{problem}

Describe in System Verilog and simulate the FIFO represented in Figure \ref{fifo}.

$\diamond$
\end{problem}

\begin{problem}\label{lb:salu}

Substitute in a RALU unit the register file with a LIFO (stack). Results a SALU in which the 16 bits used to command the register file are replaced with the 3-bit command {\tt com[2:0]}. It is required:
\begin{enumerate}
    \item prove that the SALU circuit belongs to 2-OS
    \item write the System Verilog description of SALU
    \item provide the schematic of SALU
    \item simulate SALU to prove the correctness of its behavior.
\end{enumerate}

$\diamond$
\end{problem}